\section{Founder Adaptation: de una organización con certezas a un sistema de baja probabilidad}

Después, la clase pasa al papel del fundador. Cuando McKinnon dejó Salesforce, perdió el equipo, la marca, el apoyo administrativo y una trayectoria de ascenso muy probable; las entradas del sistema emprendedor pasaron a ser un runway limitado, un buyer signal difuso y unas hipótesis que cambian todo el tiempo. Esta sección conserva esa teacher voice porque explica cómo una empresa de infraestructura, en sus primeras etapas, puede lograr avances verificables aun sin contar con evidencia completa.

\subsection{Resource Reset y daily evidence}

El fundador no puede usar la «convicción» como pretexto para rechazar los datos. Un ciclo más sólido consiste en descomponer la gran visión en evidence que pueda actualizarse cada día: quién está dispuesto a dar una entrevista, qué pain desencadena un presupuesto, si el prototype se conecta a un directory real, si el first tenant es estable, si el primer pago se repite. Cada paso reduce un tipo de incertidumbre y da fundamento a la siguiente ronda de contrataciones, de financiamiento o de decisiones de producto.

\lecturefigure{06-founder-adaptation.png}{Pasar de una empresa consolidada a un equipo emprendedor significa que los recursos y la certeza vuelven a cero; después, el sistema de ejecución se reconstruye con daily evidence.}{Subtítulos locales 06:00--12:20; concepto redibujado.}

\begin{knowledgebox}{Lectura de la figura: Belief y probability pueden coexistir}
Resource reset indica que las ventajas de la organización anterior no se trasladan de forma automática; daily evidence convierte la visión en señales de customer/prototype; low odds exige definir con claridad el runway y la stopping rule; team belief, por su parte, requiere rumbo, honestidad y avances visibles. Lo que dice el ponente, «hay que creer incluso cuando no crees», no debe entenderse como engañar al equipo, sino como la capacidad del líder de organizar la acción aun reconociendo que las probabilidades están en contra.
\end{knowledgebox}

\teachervoice{McKinnon reconoce con franqueza que su motivación para irse fue tanto la oportunidad técnica como el deseo personal de «ser su propio jefe»; lo verdaderamente difícil fue tomar una decisión en común frente a la crisis financiera, un recién nacido y los riesgos para la familia. La clase saca el founder risk del relato del héroe individual y lo devuelve a la familia, los recursos y el costo de oportunidad.}

\begin{warningbox}{No conviertas el survivorship bias en una fórmula para emprender}
El éxito posterior de Okta no demuestra que «dejar una gran empresa», «perseverar lo suficiente» o «la experiencia de la madurez» produzcan resultados por necesidad. Los principios transferibles son elegir un market acorde con las propias capacidades, obtener buyer evidence lo antes posible, limitar los riesgos irreversibles y actualizar el plan cuando aparezca nueva evidencia; el resultado sigue dependiendo del momento, la competencia, el capital y la suerte.
\end{warningbox}

\subsection{Resumen de la sección}

Founder adaptation es pasar de los inherited resources a un evidence loop. La convicción sostiene la acción continua; la probabilidad y las métricas se encargan de corregir el rumbo. Si falta cualquiera de las dos, la organización pierde el sentido de la realidad.
